%
% 1-potenzreihen.tex -- Potenzreihenentwicklung holomorpher Funktionen
%
% (c) 2025 Prof Dr Andreas Müller
%
\section{Potenzreihenansatz für Differentialgleichungen
\label{buch:differentialgleichungen:section:potenzreihen}}
\kopfrechts{Potenzreihenansatz für Differentialgleichungen}
Wir betrachten eine gewöhnliche Differentialgleichung
\begin{equation}
F(x,y,y',\dots,y^{(n)}) = 0
\label{buch:differentialgleichungen:potenzreihenansatz:eqn:dgl}
\end{equation}
der Ordnung $n$.
Wir nehmen an, dass sich die Lösung als konvergente Potenzreihe in
der Form
\begin{equation}
y(x)
=
\sum_{k=0}^\infty a_kx^k
\label{buch:differentialgleichungen:potenzreihen:eqn:ansatz}
\end{equation}
schreiben lässt.
Die Funktion $y(x)$ ist also reell analytisch und damit beliebig
oft stetig differenzierbar.
Die Exponentialfunktion oder die trigonometrischen Funktionen Sinus
und Kosinus treten als Lösungen von gewöhnlichen Differentialgleichungen
auf und haben solche konvergenten Potenzreihenentwicklungen.

Dann kann die Potenzreihe der Lösung auch als holomorphe Funktion
betrachten, indem man für $x$ auch komplexe Argumente innerhalb des
Konvergenzkreises zulässt.
Die Ableitungen $y'(x)$, $y''(x)$,\dots,$y^{(n)}(x)$ sind auch
komplexe Ableitungen.
Lässt sich auch die Funktion $F$ als eine holomorphe Funktion von
komplexen Argumenten $x$, $y$, $y'$,\dots,$y^{(n)}$ interpretieren,
dann wird
\eqref{buch:differentialgleichungen:potenzreihenansatz:eqn:dgl}
zu einer komplexen Differentialgleichung für eine gesuchte
Funktion.
Da diese holomorph sein muss, ist es auch naheliegend, sie als
Potenzreihe zu schreiben.

Der angesprochene Fall, dass in $F$ komplexe Argumente eingesetzt
werden können, ist zum Beispiel bei einer linearen Differentialgleichung
\[
y^{(n)}(x)
+
a_{n-1}(x) y^{(n-1)}(x)
+
\dots
+
a_2(x) y''(x)
+
a_1(x) y'(x)
+
a_0(x)
=
0
\]
gegeben, sofern die Koeffizientenfunktionen $a_k(x)$ holomorphe Funktionen
eines komplexen Argumentes $a_k(x)$ sind.
In den nachfolgenden Abschnitten werden zum Beispiele für lineare
Differentialgleichungen mit konstanten Koeffizienten gezeigt.
Die fuchsche Differentialgleichung, die in
Abschnitt~\ref{buch:differentialgleichungen:section:fuchs} behandelt wird,
ist eine lineare Differentialgleichung mit
rationalen oder analytische Koeffizientenfunktionen.

%
% Die Exponentialfunktion
%
\subsection{Die Exponentialfunktion}
Wir betrachten die Differentialgleichung
\begin{equation}
y' = -ky
\label{buch:differentialgleichungen:potenzreihen:eqn:expdgl}
\end{equation}
und setzen den
Potenzreihenansatz~\eqref{buch:differentialgleichungen:potenzreihen:eqn:ansatz}
in die Differentialgleichung ein.
Die Ableitung der  
Potenzreihen~\eqref{buch:differentialgleichungen:potenzreihen:eqn:ansatz}
ist
\[
y'(x)
=
\sum_{k=1}^\infty ka_kx^{k-1}.
\]
Einsetzen in die Differentialgleichung
\eqref{buch:differentialgleichungen:potenzreihen:eqn:expdgl}
ergibt die Gleichung
\begin{align*}
\sum_{l=1}^\infty a_llx^{l-1}
&=
-k
\sum_{l=0}\infty a_lx^l.
\intertext{Um diese besser vergleichen zu können, ersetzen wir $l-1$ durch
$l$, oder $l$ durch $l+1$ auf der linken Seite und erhalten}
\sum_{l=0}^\infty a_{l+1} (l+1)x^{l}
&=
-k
\sum_{l=0}^\infty a_lx^l.
\end{align*}
Durch Vergleich der Koeffizienten der beiden Potenzreihen können wir
schliessen, dass
\begin{equation}
a_{l+1}
=
-k
\frac{1}{l+1} a_l
\label{buch:differentialgleichungen:potenzreihen:eqn:exprek}
\end{equation}
ist für alle $l\ge 0$.
\eqref{buch:differentialgleichungen:potenzreihen:eqn:exprek}
ist eine Rekursionsgleichung, die ermöglicht, die Koefizienten $a_k$
mit $k>0$ auf $a_0$ zurückzuführen.
Löst man die Gleichung iterativ auf, erhält man die Werte
\begin{equation}
\begin{aligned}
a_1
&=
-k\frac{1}{0+1}a_0 = -ka_0 
\\
a_2
&=
-k\frac{1}{1+1}a_1 = -\frac{k}{2}a_1 = \phantom{-}\frac{k^2}{2}a_0
\\
a_3
&=
-k\frac{1}{2+1}a_2 = -\frac{k}{3}a_2 = -\frac{k^3}{3!}a_0
\\
a_4
&=
-k\frac{1}{3+1}a_3 = -\frac{k}{4}a_3 = \phantom{-}\frac{k^4}{4!}a_0
\\
a_5
&=
-k\frac{1}{4+1}a_4 = -\frac{k}{5}a_4 = -\frac{k^5}{5!}a_0
\\
&\hspace*{1mm}\vdots
\\
a_l
&=
(-1)^l \frac{k^l}{l!}a_0.
\end{aligned}
\label{buch:differentialgleichungen:potenzreihen:eqn:rekursion}
\end{equation}
Der Koeffizient $a_0$ wird durch die Differentialgleichung oder
durch \eqref{buch:differentialgleichungen:potenzreihen:eqn:expdgl}
nicht festgelegt.

Aus \eqref{buch:differentialgleichungen:potenzreihen:eqn:rekursion}
ergibt sich als Lösung der Differentialgleichung als die Potenzreihe
\[
y(x)
=
\sum_{l=0}^\infty a_lx^l
=
a_0
\sum_{l=0}^\infty
(-1)^l
\frac{k^l}{l!}x^l
=
a_0
\sum_{l=0}^\infty
\frac{(-kx)^l}{l!}
=
a_0e^{-kx}.
\]
Tatsächlich ist die Ableitung dieser Funktion
\[
y'(x)
=
-ka_0e^{-kx}
=
-ky(x),
\]
die Exponentialfunktion löst also die Differentialgleichung.

Da 
\eqref{buch:differentialgleichungen:potenzreihen:eqn:expdgl}
eine Differentialgleichung erster Ordnung ist, erwartet man 
genau eine frei wählbare Integrationskonstante.
Der Koeffizient $a_0$ ist diese Konstante.
Sie kann aus einer Anfangsbedingung $y(0)=a_0$ bestimmt werden.

%
% Schwingungen
%
\subsection{Schwingungen}
Die Schwingungsdifferentialgleichung ist eine Differentialgleichung
der Form
\begin{equation}
y''
=
-\omega^2\,y.
\label{buch:differentialgleichungen:potenzreihen:eqn:schwingungsdgl}
\end{equation}
Auch in diesem Fall sind die Koeffizienten konstant und somit
analytisch und es ist gerechtfertigt, eine analytische Lösung
zu erwarten und mit einem Potenzreihenansatz zu suchen.

%
% Reell analytische Lösung
%
\subsubsection{Reell analytische Lösung}
Die zweite Ableitung des
Potenzreihenansatzes~\eqref{buch:differentialgleichungen:potenzreihen:eqn:ansatz}
ist
\[
y''(x)
=
\sum_{l=2}^\infty
a_l l(l-1) x^{l-2}
.
\]
Setzen wir dies in die Differentialgleichung
\eqref{buch:differentialgleichungen:potenzreihen:eqn:schwingungsdgl}
ein, ergibt sich
\begin{align*}
\sum_{l=2}^\infty
a_l l(l-1) x^{l-2}
&=
-\omega^2
\sum_{l=0}^\infty
a_lx^l
\intertext{Auch in diesem Fall ist es notwendig, auf der rechten
Seite $l+2$ durch $l$ zu ersetzen, was auf die Gleichung}
\sum_{l=0}^\infty
a_{l+2} (l+1)(l+2) x^{l}
&=
-\omega^2
\sum_{l=0}^\infty
a_lx^l
\end{align*}
Koeffizientenvergleich ergibt die Rekursionsformel
\begin{align*}
(l+1)(l+2)
a_{l+2} 
&=
-\omega^2
a_l
\intertext{für $l>0$ oder}
a_{l+2}
&=
-\frac{\omega^2}{(l+1)(l+2)} a_{l}.
\end{align*}
Man beachte, dass die Rekursionsformel nur entweder gerade oder ungerade
indizierte Koeffizienten untereinander in Beziehung setzt.
Die gerade indizierten Koeffizienten müssen sich daher alle durch
$a_0$ ausdrücken lassen, während die ungerade indizierten Koeffizienten
von $a_1$ abhängen.
Die Lösungsfunktion ist also erst durch die Angabe von $a_0=y(0)$ und
$a_1=y'(0)$ eindeutig festgelegt.

%
% Lösung der Rekursionsgleichung
%
\subsubsection{Lösung der Rekursionsgleichung}
Diese Rekursionsformeln ermöglichen, die Koeffizienten $a_l$ mit $l\ge 2$
durch $a_0$ und $a_1$ als
\begin{align*}
a_2
&=
-\omega^2 \frac{1}{1\cdot 2} a_0
\\
a_3
&=
-\omega^2 \frac{1}{2\cdot 3} a_1
%=
%-\omega^2 \frac{1}{3!} a_1
\\
a_4
&=
-\omega^2 \frac{1}{3\cdot 4} a_2
=
\phantom{-}\omega^4 \frac{1}{4!} a_0
\\
a_5
&=
-\omega^2 \frac{1}{4\cdot 5} a_3
=
\phantom{-}\omega^4 \frac{1}{5!} a_1
\\
a_6
&=
-\omega^2 \frac{1}{5\cdot 6} a_4
=
-\omega^6 \frac{1}{6!} a_0
\\
a_7
&=
-\omega^2 \frac{1}{6\cdot 7} a_5
=
-\omega^6 \frac{1}{7!} a_1
\intertext{auszudrücken oder allgemein}
a_{2n}   &= (-1)^n \omega^{2n} \frac{1}{(2n)!}   a_0
\\
a_{2n+1} &= (-1)^n \omega^{2n} \frac{1}{(2n+1)!} a_1.
\end{align*}
Daraus kann die Lösung als
\begin{align}
y(x)
&=
a_0
\sum_{n=0}^\infty
(-1)^n\omega^{2n}\frac{1}{(2n)!} x^{2n}
+
a_1
\sum_{n=0}^\infty
(-1)^n\omega^{2n}\frac{1}{(2n+1)!} x^{2n+1}
\notag
\\
&=
a_0
\sum_{n=0}^\infty
\frac{(-1)^n}{(2n)!} (\omega x)^{2n}
+
\frac{a_1}{\omega}
\sum_{n=0}^\infty
\frac{(-1)^n}{(2n+1)!} (\omega x)^{2n+1}
\notag
\\
&=
a_0\cos(\omega x)
+
\frac{a_1}{\omega}\sin(\omega x)
\label{buch:differentialgleichungen:potenzreihen:eqn:schwloesung}
\end{align}
ableiten.
Die Lösungen sind also Linearkombinationen von Sinus-
bzw.~Kosinus-Schwin\-gun\-gen.

Auch in diesem Fall ist die zweite Ableitung der Lösung
\eqref{buch:differentialgleichungen:potenzreihen:eqn:schwloesung}
\[
y''(x)
=
-\omega^2 a_0\cos(\omega x)
-
\omega^2 \frac{a_1}{\omega} \sin(\omega x)
=
-\omega^2 y(x),
\]
und erfüllt damit wie erwartet die Schwingungsdifferentialgleichung.

%
% Komplexe Lösungen
%
\subsubsection{Komplexe Lösungen}
Die Differentialgleichung 
\eqref{buch:differentialgleichungen:potenzreihen:eqn:schwingungsdgl}
kann man auch in der Form
\begin{equation}
\biggl(\frac{d}{dx}-i\omega\biggr)
\biggl(\frac{d}{dx}+i\omega\biggr)
y
=
0
\label{buch:differentialgleichungen:potenzreihen:eqn:faktschwing}
\end{equation}
schreiben.
Lösungen der Differentialgleichungen erster Ordnung
\begin{equation}
\biggl(\frac{d}{dx} - i \omega\biggr)y = 0
\qquad\text{und}\qquad
\biggl(\frac{d}{dx} + i \omega\biggr)y = 0
\label{buch:differentialgleichungen:potenzreihen:eqn:einzelschwing}
\end{equation}
ergeben linear unabhängige Lösungen der Differentialgleichung
\eqref{buch:differentialgleichungen:potenzreihen:eqn:faktschwing}.
Die Differentialgleichungen
\eqref{buch:differentialgleichungen:potenzreihen:eqn:einzelschwing}
sind komplexe Differentialgleichungen der Form
\eqref{buch:differentialgleichungen:potenzreihen:eqn:expdgl}
und können daher mit der Exponentialreihe gelöst werden.
Die Lösungen von 
\eqref{buch:differentialgleichungen:potenzreihen:eqn:einzelschwing}
\begin{equation}
y_+(x) = e^{i\omega x}
\qquad\text{und}\qquad
y_-(x) = e^{-i\omega x}
\end{equation}
können zu den allgemeinen Lösungen der Schwingungsdifferentialgleichungen
linear kombiniert werden.
Die trigonometrischen Funktionen als Lösungen entstehen als die
Linearkombinationen
\[
\cos \omega x
=
\frac{y_+(x)+y_-(x)}{2}
\qquad\text{und}\qquad
\sin \omega x
=
\frac{y_+(x)-y_-(x)}{2i}.
\]

%
% Die Wurzelfunktion
%
\subsection{Die Wurzelfunktion}
Die Wurzelfunktion ist eine Lösung der algebraischen Gleichung $y^2=x$.
Durch Ableiten nach $x$ wird daraus die nichtlineare Differentialgleichung
\begin{equation}
2yy'=1.
\label{buch:differentialgleichungen:potenzreihen:eqn:wurzeldgl}
\end{equation}
Mit der Funktion
\[
F(x,y,y')
=
2yy'-1
\]
kann die Differentialgleichung
\eqref{buch:differentialgleichungen:potenzreihen:eqn:wurzeldgl}
in der Form
\eqref{buch:differentialgleichungen:potenzreihenansatz:eqn:dgl}
geschrieben werden.
Die Funktion $F$ ist eine Polynom in den Variablen $y$ und $y'$ und
ist damit insbesondere auch für komplexe Argumente definiert.
\eqref{buch:differentialgleichungen:potenzreihen:eqn:wurzeldgl}
ist eine komplexe Differentialgleichung für die Wurzelfunktion
und sollte daher eine komplexe Potenzreihe für die Wurzelfunktion
zu liefern in der Lage sein.

%
% Potenzreihenlösung für die Wurzelfunktion
%
\subsubsection{Potenzreihenlösung für die Wurzelfunktion}
Wir schreiben die Lösung als
\[
y(x)
=
\sum_{k=0}^\infty a_kx^k
\qquad\Rightarrow\qquad
y'(x)
=
\sum_{k=1}^\infty ka_kx^{k-1}
=
\sum_{l=0}^\infty (l+1)a_{l+1}x^l
\]
und setzen in die Differentialgleichung
\eqref{buch:differentialgleichungen:potenzreihen:eqn:wurzeldgl}
ein.
Es entsteht die Gleichung
\begin{align}
\frac12
&=
\sum_{k=0}^\infty \sum_{l=0}^\infty a_k(l+1)a_{l+1}x^{k+l}
=
\sum_{n=0}^\infty
\biggl(
\sum_{l=0}^n (l+1)a_{n-l}a_{l+1}
\biggr)
x^n.
\end{align}
Durch Koeffizientenvergleich können wir jetzt die Koeffizienten $a_k$
bestimmen.


Der Koeffizient für $n=0$ ist $a_0a_1=\frac12$ oder $a_1=\frac{1}{2a_0}$.
Die Koeffizienten von $x^n$ mit $n>0$ müssen alle verschwinden, es muss
also
\begin{align*}
0
&=
\sum_{l=0}^n (l+1)a_{n-l}a_{l+1}
&&\Rightarrow&
(n+1)a_0a_{n+1}
&=
-
\sum_{l=0}^{n-1} (l+1)a_{n-l}a_{l+1}
\\
&
&&\Rightarrow&
a_{n+1}
&=
\frac{1}{(n+1)a_0}
\sum_{l=0}^{n-1} (l+1)a_{n-l}a_{l+1}.
\end{align*}
Dies ist eine Rekursionsformel, die den Koeffizienten $a_{n+1}$ aus
den Koeffizienten $a_0,\dots,a_n$ bestimmt.
Alle Koeffizienten hängen also vom Wert $a_0$ ab.

%
% Der Spezialfall a_0=1
%
\subsubsection{Der Spezialfall $a_0=1$}
Um die Rechnung zu vereinfachen, berechnen wir die Koeffizienten nur für 
den Fall $a_0=1$, also für den Fall einer Wurzelfunktion, die für $x=0$
den Wert $1$ ergibt.
Um die Koeffizienten besser auseinanderzuhalten, bezeichnen wir die
Koeffizienten für diesen Spezialfall mit $b_k$.
\input{chapters/050-differentialgleichungen/fig/fig-wurzel.tex}%
Wir erwarten also eine Potenzreihenentwicklung für die Funktion
$y(x)=\!\sqrt{1+x}$.
Die Rechnung ergibt
\begin{align*}
n&=0: &&&
1 &=
b_0{\color{darkred}b_1}
\\
&&&\Rightarrow&
b_1
&=
\frac{1}{2b_0}
&&=
\frac12
\\
n&=1: &&&
0 &=
b_1^2+2b_0{\color{darkred}b_2}
\\
&&&\Rightarrow&
{\color{darkred}b_2}
&= 
-\frac{b_1^2}{2b_0}
&&=
-\frac{1}{8}
\\
n&=2: &&&
0 &=
b_2b_1+2b_1b_2+3b_0{\color{darkred}b_3}
\\
&&&\Rightarrow&
{\color{darkred}b_3}
&=
-\frac{1}{3b_0}(3b_1b_2)
&&=
\frac{1}{16}
\\
n&=3: &&&
0 &=
b_3b_1+2b_2^2+3b_1b_3+4b_0{\color{darkred}b_4}
\\
&&&\Rightarrow&
{\color{darkred}b_4}
&=
-
\frac{1}{4b_0}(4b_3b_1+2b_2^2)
&&=
-\frac{5}{128}
\\
n&=4: &&&
0 &=
b_4b_1+2b_3b_2+3b_2b_3+4b_1b_4+5b_0{\color{darkred}b_5}
\\
&&&\Rightarrow&
{\color{darkred}b_5}
&=
-\frac{1}{5b_0}(5b_1b_4+5b_2b_3)
&&=
\frac{7}{256}
\end{align*}
Die Lösung hat also die Potenzreihe
\begin{align}
\!\sqrt{1+x}
&=
\frac12
-\frac{1}{8}x
+\frac{1}{16}x^2
-\frac{5}{128}x^3
+\frac{7}{256}x^4
-\dots
\label{buch:differentialgleichungen:potenzreihen:wurzelapprox}
\end{align}
Die Abbildung~\ref{buch:differentialgleichungen:fig:wurzel}
zeigt die exakte Lösung $\!\sqrt{1+x}$ im Vergleich zur Approximation
\eqref{buch:differentialgleichungen:potenzreihen:wurzelapprox}.

%
% Der allgemeine Fall
%
\subsubsection{Der allgemeine Fall}
Die Differentialgleichung
\eqref{buch:differentialgleichungen:potenzreihen:eqn:wurzeldgl}
hat als allgemeine Lösung die Funktion $y(x)=\!\sqrt{x-C}$.
Die Wahl $a_0=1$ entspricht daher dem Wert $C=-1$.
Für beliebiges $a_0$ findet man
\[
y(x)
=
\!\sqrt{x+a_0^2}
=
\frac{1}{a_0}
\sqrt{1+\frac{x\mathstrut}{a_0^2\mathstrut}}
=
\frac{1}{a_0}
\sum_{k=0}^\infty b_k\biggl(\frac{x\mathstrut}{{a_0}^2\mathstrut}\biggr)^k
\]
als Lösung.
Daraus können wir ablesen, dass
\[
a_k
=
\frac{b_k}{a_0^{2k+1}}
\]
ist.

%
% Die binomische Reihe
%
\subsubsection{Die binomische Reihe}
Für die Funktion $f(x) = (1+x)^\alpha$ ist auch die binomische Formel
\begin{align*}
(1+x)^\alpha
&=
\sum_{k=0}^\infty \binom{\alpha}{k}x^k
=
\sum_{k=0}^\infty \frac{\alpha(\alpha-1)\cdot\ldots\cdot (\alpha-k+1)}{k!}x^k
\end{align*}
bekannt
Im Spezialfall $\alpha=\frac12$ kann man die Binomialkoeffizienten
explizit ausrechnen.
Dazu erweitern wir den Bruch
\begin{align}
\binom{\frac12}{k}
&=
\frac{\frac12\cdot(-\frac12)\cdot(-\frac32)\cdot\ldots\cdot(-\frac{2k-3}{2})}{k!}
=
(-1)^{k-1}
\frac{1\cdot 3 \cdot\ldots\cdot(2k-3)}{2^k\cdot k!}
\intertext{mit $2^k\cdot k!=2\cdot 4 \cdot \ldots\cdot (2k)$ und erhalten}
&=
(-1)^{k-1}
\frac{1\cdot 2\cdot 3 \cdot 4 \cdot\ldots\cdot(2k-3)\cdot(2k-2)\cdot 2k}{(2^k\cdot k!)^2}
\notag
\\
&=
(-1)^{k-1}
\frac{(2k)!}{2^{2k}\cdot (k!)^2\cdot(2k-1)}
=
(-1)^{k-1}
\frac{1}{2^{2k}}
\binom{2k}{k}.
\label{buch:differentialgleichungen:potenzreihen:eqn:koeffizient}
\end{align}
Tatsächlich bekommt
\eqref{buch:differentialgleichungen:potenzreihen:eqn:koeffizient}
für die ersten paar $k$ die Werte der Koeffizienten der
Tabelle~\ref{buch:differentialgleichungen:tab:wurzel}.
\input{chapters/050-differentialgleichungen/fig/tab-wurzel.tex}%
Sie stimmen mit den Koeffizienten in
\eqref{buch:differentialgleichungen:potenzreihen:wurzelapprox}
überein.

